\documentclass{ltjbook}
\usepackage{docmute} 
\usepackage{../../myStyle} 

\begin{document}
\chapter{JASPによるベイズ推定の実践}
% \item[JASPの導入] JASPのインストールと起動，その基本的な特徴の解説など導入する演習を行う。各自のPCにJASP環境を導入し，サンプルコードを開いて記述統計量やプロットを描く。現段階\footnote{2021.12.01現在}ではJASPを用いたテキストや解説書は少ないが，標準的なアプリケーションの導入で対応できるため，初学者にもその障壁は低いと考えられる。
% \item[JASPによる回帰分析] JASPの特徴の一つは，従来法とベイズ法が同じメニューから並列的に利用できることである。ここでは回帰分析を例に，最尤法とベイズ法で結果の表示の違いや解釈の方法について見ていく。
% \item[JASPによる平均の比較]
\end{document}
